% =====================================================================
%  Chương 3 — Hạng mục 3: hiện thực và huấn luyện bằng PyTorch
%  Mã nguồn: a06/models_torch.py, a06/train_torch.py, a06/bench.py
%  Notebook: 03_pytorch_amzn.ipynb, 04_pytorch_kkbox.ipynb
% =====================================================================
\chapter{Hiện thực và huấn luyện bằng PyTorch}
\label{ch:pytorch}

Chương này dựng bốn mô hình Simple RNN, LSTM, GRU và BiLSTM bằng PyTorch~\cite{paszke2019}, huấn luyện trên cả hai
tập của Chương~\ref{ch:dulieu} và đánh giá trên tập test. Bốn mô hình chỉ khác nhau ở tầng hồi quy; phần còn lại
(đầu dự báo, tối ưu, dữ liệu) giống hệt, để chênh lệch kết quả chỉ đến từ loại tế bào.

\section{Một lớp cho bốn mô hình}

Mô hình nhận tensor $(B, T, 8)$, cho qua tầng hồi quy 64 đơn vị, lấy vector tóm tắt $h_T$ (hoặc
$[\overrightarrow{h}_T; \overleftarrow{h}_1]$ với BiLSTM) rồi đưa qua đầu dự báo
$\mathrm{Linear}(\cdot, 32) \to \mathrm{ReLU} \to \mathrm{Linear}(32, 1)$. Đầu ra là một số: return đã chuẩn hoá với
AMZN, logit churn với KKBox. Cả bốn mô hình dùng chung một lớp; tham số \texttt{cell} chọn tầng hồi quy.

\maNguon{models_torch__PyTorchRNN}{Lớp chung cho Simple RNN, LSTM, GRU và BiLSTM}

\texttt{nn.LSTM} và \texttt{nn.GRU} gọi thẳng nhân cuDNN khi chạy trên GPU, nên vòng lặp theo thời gian không nằm
trong Python. Với BiLSTM, PyTorch trả về đầu ra hai chiều nối nhau ở mỗi bước; lớp trên cắt đúng
$\overrightarrow{h}_T$ ở bước cuối và $\overleftarrow{h}_1$ ở bước đầu, đúng như cách nối đã kiểm chứng ở
Chương~\ref{ch:lythuyet}.

\section{Vòng lặp huấn luyện}

Cấu hình giống nhau cho mọi mô hình: AdamW~\cite{loshchilov2019} với tốc độ học $10^{-3}$ và weight decay $10^{-4}$, cắt chuẩn gradient về 1 (Chương~\ref{ch:lythuyet}), dừng sớm khi val loss không
giảm sau một số epoch, và giữ lại trọng số của epoch có val loss nhỏ nhất. Bảng~\ref{tab:cfg} ghi phần khác nhau giữa hai
tập.

\begin{table}[H]
  \centering\small
  \caption{Cấu hình huấn luyện}
  \label{tab:cfg}
  \begin{tabular}{lll}
    \toprule
    & AMZN (hồi quy) & KKBox (phân loại) \\
    \midrule
    Hàm mất mát & MSE trên return chuẩn hoá & BCE trên logit, \texttt{pos\_weight} = số âm / số dương \\
    Lô, số epoch tối đa & 64, 60 & 512, 20 \\
    Dừng sớm & 8 epoch không cải thiện & 4 epoch không cải thiện \\
    Thiết bị & chọn bằng đo (Mục~\ref{sec:thietbi}) & GPU \\
    Seed & 42 & 42 \\
    \bottomrule
  \end{tabular}
\end{table}

Toàn bộ dữ liệu được chuyển lên thiết bị một lần dưới dạng tensor; mỗi epoch chỉ xáo chỉ số rồi cắt lô, không dùng
\texttt{DataLoader}. Với tập KKBox khoảng 100 MB, cách này bỏ được chi phí chép dữ liệu từ RAM sang GPU ở mỗi lô.

\maNguon{train_torch__fit}{Vòng lặp huấn luyện chung: AdamW, clip gradient, dừng sớm theo val loss}

Với KKBox, nhãn dương chỉ chiếm khoảng 9\,\%. Không cân bằng lại thì mô hình đạt loss thấp bằng cách đoán ``ở lại'' cho
tất cả. \texttt{pos\_weight} nhân mất mát của mẫu dương lên khoảng 10 lần, tương đương cân bằng hai lớp trong hàm mất mát.

\maNguon{train_torch__make_loss}{Hàm mất mát cho hai loại bài toán}

\subsection{Chọn thiết bị bằng đo}
\label{sec:thietbi}

Các bài trước cho thấy GPU không phải lúc nào cũng nhanh hơn với mô hình nhỏ và tập nhỏ. Vì vậy notebook 03 đo một
epoch LSTM trên cả hai thiết bị (sau một lượt khởi động) rồi chọn bên nhanh hơn. Kết quả: CPU
\SL{torch_amzn/probe/cpu} s/epoch, GPU \SL{torch_amzn/probe/cuda} s/epoch, nên AMZN cũng huấn luyện trên GPU. Nhân cuDNN
cho LSTM chạy cả 30 bước trong một lần gọi, còn CPU phải lặp từng bước.

\maNguon{train_torch__pick_device}{Chọn CPU hay GPU bằng cách đo một epoch}

\section{Kết quả trên AMZN}

\begin{figure}[H]
  \centering
  \hinh{f3_loss_amzn}
  \caption{Đường loss của bốn mô hình PyTorch trên AMZN (liền: train, đứt: val)}
  \label{fig:loss-amzn}
\end{figure}

Hình~\ref{fig:loss-amzn} cho thấy train loss giảm chậm nhưng đều, còn val loss gần như không đổi ngay từ những epoch
đầu. Val loss thấp hơn train loss nhiều vì nhãn được chuẩn hoá theo độ lệch chuẩn của train, trong khi return giai đoạn
val biến động ít hơn hẳn (Chương~\ref{ch:dulieu}); so hai đường với nhau thì không có ý nghĩa, chỉ xu hướng của từng
đường là có. BiLSTM có val loss nhỏ nhất ở epoch \SL{torch_amzn/models/bilstm/best_epoch}, tức là mọi epoch sau đó chỉ
làm mô hình khớp thêm nhiễu của tập train. Các mô hình còn lại đạt val loss nhỏ nhất ở epoch
\SL{torch_amzn/models/lstm/best_epoch} đến \SL{torch_amzn/models/rnn/best_epoch}.

\begin{table}[H]
  \centering\footnotesize
  \caption{Bốn mô hình PyTorch trên tập test AMZN (\SL{torch_amzn/test_start} đến \SL{torch_amzn/test_end})}
  \label{tab:torch-amzn}
  \input{generated/tab_torch_amzn}
\end{table}

Bảng~\ref{tab:torch-amzn} là kết quả trung thực nhất của bài, và nó không đẹp. RMSE của bốn mô hình nằm trong khoảng
\SL{torch_amzn/models/bilstm/rmse} đến \SL{torch_amzn/models/gru/rmse} USD, trong khi mốc ngây thơ ``giá ngày mai bằng
giá hôm nay'' đạt \SL{torch_amzn/naive/rmse} USD. Không mô hình nào thắng mốc ngây thơ. Cột $R^2$ trên giá cho
\SL{torch_amzn/models/lstm/r2} với mọi mô hình và với cả mốc ngây thơ, nên con số rất cao này không nói gì: giá hôm
nay đã giải thích gần hết phương sai của giá ngày mai. Cột $R^2$ trên return mới đo khả năng dự báo thật, và nó âm với
cả bốn mô hình. Tỉ lệ dự báo đúng chiều tăng hay giảm dao động quanh 50\,\%, từ
\SL{torch_amzn/models/rnn/dir_acc}\,\% đến \SL{torch_amzn/models/lstm/dir_acc}\,\%.

\begin{figure}[H]
  \centering
  \hinh{f3_pred_amzn}
  \caption{Giá thật, mốc ngây thơ và dự báo của bốn mô hình trên các phiên cuối của tập test}
  \label{fig:pred-amzn}
\end{figure}

Hình~\ref{fig:pred-amzn} cho thấy nhìn bằng mắt thì dự báo ``bám'' giá thật rất sát. Phóng to 30 phiên cuối (ô b) thì
thấy rõ lý do: đường dự báo gần như trùng mốc ngây thơ, tức là đường giá thật dịch sang phải một phiên. Mô hình học được
rằng return ngày mai có trung bình gần 0, và dự báo gần 0 là lựa chọn có sai số bình phương nhỏ nhất khi return không có
tự tương quan (Chương~\ref{ch:dulieu}). Chương~\ref{ch:keras} bàn tiếp giới hạn này.

\section{Kết quả trên KKBox}

\begin{figure}[H]
  \centering
  \hinh{f3_loss_kkbox}
  \caption{Đường loss của bốn mô hình PyTorch trên KKBox (liền: train, đứt: val)}
  \label{fig:loss-kkbox}
\end{figure}

KKBox là bài toán khác hẳn. Train loss và val loss giảm cùng nhau trong nhiều epoch (Hình~\ref{fig:loss-kkbox}); LSTM
đạt val loss tốt nhất ở epoch \SL{torch_kkbox/models/lstm/best_epoch}/\SL{torch_kkbox/models/lstm/epochs_run}. Mô hình
học được tín hiệu thật từ chuỗi hành vi.

Ngưỡng phân loại không cố định ở 0,5. Với mỗi mô hình, notebook chọn ngưỡng làm F1 lớn nhất trên \emph{validation},
rồi áp nguyên ngưỡng đó cho test. Chọn ngưỡng trên test sẽ làm kết quả test đẹp hơn thực tế.

\maNguon{bench__best_threshold}{Chọn ngưỡng theo F1 trên validation}

\begin{table}[H]
  \centering\footnotesize
  \caption{Bốn mô hình PyTorch trên tập test KKBox (\SL{kkbox/n_test} người dùng)}
  \label{tab:torch-kkbox}
  \input{generated/tab_torch_kkbox}
\end{table}

Ba mô hình có cổng (LSTM, GRU, BiLSTM) đạt ROC-AUC từ \SL{torch_kkbox/models/gru/roc_auc} đến
\SL{torch_kkbox/models/lstm/roc_auc}, cao hơn Simple RNN (\SL{torch_kkbox/models/rnn/roc_auc}). Khoảng cách không lớn
vì chuỗi chỉ dài 31 bước, chưa đủ dài để gradient của Simple RNN triệt tiêu hoàn toàn. LSTM có F1 cao nhất,
\SL{torch_kkbox/models/lstm/f1}\,\%, với Precision \SL{torch_kkbox/models/lstm/precision}\,\% và Recall
\SL{torch_kkbox/models/lstm/recall}\,\%: cứ ba người bị báo churn thì khoảng một người churn thật, và mô hình tìm ra
khoảng một phần ba số người churn.

\begin{figure}[H]
  \centering
  \hinh{f3_conf_kkbox}
  \caption{Ma trận nhầm lẫn trên test KKBox; màu ô theo tỉ lệ trong hàng}
  \label{fig:conf-kkbox}
\end{figure}

\begin{figure}[H]
  \centering
  \hinh{f3_roc_pr}
  \caption{Đường ROC và Precision--Recall trên test KKBox; vạch đứt là mô hình đoán ngẫu nhiên}
  \label{fig:roc-pr}
\end{figure}

Các con số này thấp so với bảng xếp hạng của cuộc thi, và có lý do. Các đội mạnh nhất dùng thêm lịch sử giao dịch (phương
thức thanh toán, tự động gia hạn, ngày hết hạn), là những đặc trưng dự báo churn trực tiếp. Bài này chỉ dùng log nghe
nhạc, vì trọng tâm của bài là mô hình hồi quy trên dữ liệu chuỗi. Hình~\ref{fig:roc-pr} cho thấy phần tín hiệu mà chuỗi hành vi
mang lại là có thật: đường PR nằm cao hơn nhiều so với mốc ngẫu nhiên \SL{torch_kkbox/pos_rate_test}\,\%, nhưng
ROC-AUC dừng ở khoảng 0,68. Phần tín hiệu còn lại không nằm trong việc người dùng nghe nhạc thế nào.
